\documentclass[12pt]{article}
\usepackage[utf8]{inputenc}
\usepackage[T1]{fontenc}
\usepackage{hyperref}
\title{Tiéfo Language Resources for K-12 Curriculum Development}
\author{Compiled by Claude Sonnet 5 (claude-sonnet-5) for Andr\'as Kornai}
\date{}
\begin{document}
\maketitle

\section{Introduction}
Tiéfo (ISO 639-3: tiq; also Tyefo, Kiefo) is an endangered language of Burkina Faso whose classification is uncertain (possibly peripheral Gur). Two mutually unintelligible varieties are reported: Tiefo-D (Daramandougou, Comoé) and Tiefo-N (Noumoudara and Nyafogo, Houet), with Nyafogo nearly extinct. About 1,000 people are reported to speak it as a mother tongue, and Jula is dominant in the region.

\section{Primary Educational Resources}
No primers, textbooks or school materials were found in search results.

\section{Linguistic and Reference Materials}
\begin{itemize}
\item Wikipedia, Tiefo language (overview; mentions Winkelmann's 1998 dissertation, a short Tiefo-N grammar from 2017, and fieldwork on Tiefo-D begun in 2017 by Heath and Ouattara): \url{https://en.wikipedia.org/wiki/Tiefo_language}
\item Verbix Endangered Languages page for Tiefo: \url{https://docs.verbix.com/EndangeredLanguages/Tiefo}
\end{itemize}

\section{Religious and Community Resources}
None found in search results.

\section{Digital and Audio Resources}
The Endangered Languages Project entry for Tiéfo is a catalogue record; the search did not confirm any audio or corpus content there: \url{https://endangeredlanguages.com/lang/3818}

\section{Curriculum Development Context}
The language is not known to be taught in schools, and speakers are mainly older adults. The only substantial scholarship found is academic documentation. Any curriculum would first require community consultation, a choice between the two varieties, and further documentation.
\end{document}